\documentclass[a4paper]{article}
% Shared preamble for the labs (compile from Labs/ with XeLaTeX or tectonic).
% Same fonts and header as the Series sheets, plus colored boxes and tables.

\usepackage{enumitem}
\usepackage[margin=2.2cm]{geometry}
\usepackage{amsmath,amsfonts,amssymb}
\usepackage[charter]{mathdesign}
\usepackage{array}
\usepackage{colortbl}
\usepackage{fancyhdr}
\usepackage[skins,breakable]{tcolorbox}
\usepackage{polyglossia}

%------------------------------------------------
% Languages and fonts
%------------------------------------------------
\setdefaultlanguage[calendar=gregorian, numerals=maghrib]{arabic}
\setotherlanguage{english}

\IfFontExistsTF{Amiri}{%
    \newfontfamily\arabicfont[Script=Arabic, Scale=1.15]{Amiri}
}{%
    \newfontfamily\arabicfont[
        Script=Arabic, Scale=1.15,
        Path=../fonts/, Extension=.ttf,
        UprightFont=*-Regular, BoldFont=*-Bold,
        ItalicFont=*-Italic, BoldItalicFont=*-BoldItalic
    ]{Amiri}
}

\IfFontExistsTF{AlBattar}{%
    \newfontfamily\arabicfontsf[Script=Arabic, Scale=1.2]{AlBattar}
    \newfontfamily\khatt[Script=Arabic, Scale=1.2]{AlBattar}
}{%
    \newfontfamily\arabicfontsf[
        Script=Arabic, Scale=1.2, Path=../fonts/,
        BoldFont=ae_AlBattar.ttf
    ]{ae_AlBattar.ttf}
    \newfontfamily\khatt[
        Script=Arabic, Scale=1.2, Path=../fonts/,
        BoldFont=ae_AlBattar.ttf
    ]{ae_AlBattar.ttf}
}

%------------------------------------------------
% Colors (same palette as the slides)
%------------------------------------------------
\definecolor{brandTeal}{HTML}{00B4D8}
\definecolor{brandOrange}{HTML}{FF7043}
\definecolor{brandDark}{HTML}{263238}
\definecolor{brandMuted}{HTML}{6C757D}
\colorlet{tealDeep}{brandTeal!70!black}

%------------------------------------------------
% Page setup
%------------------------------------------------
\renewcommand{\footrulewidth}{1pt}
\renewcommand{\headrulewidth}{1pt}
\rhead{\khatt المدرسة العليا للأساتذة ببوزريعة}
\lhead{\khatt ملحقة بومرداس}
\cfoot{\thepage}
\lfoot{\khatt الأستاذ: د. بن عبد الرزاق يوسف}
\parindent 0pt
\setlength{\headsep}{15pt}
\setlength{\headheight}{16pt}
\setlength{\parskip}{4pt}
\pagestyle{fancy}

\setlist[enumerate,1]{label=\arabic*), itemsep=4pt, topsep=3pt}
\setlist[enumerate,2]{label=\alph*), itemsep=2pt, topsep=2pt}
\setlist[itemize]{itemsep=2pt, topsep=2pt}

%------------------------------------------------
% Tables
%------------------------------------------------
\renewcommand{\arraystretch}{1.4}
\setlength{\arrayrulewidth}{0.6pt}
\arrayrulecolor{brandMuted!80}
% Header cell
\newcommand{\ras}[1]{\cellcolor{brandTeal!20}\textbf{#1}}
% Highlighted result cell
\newcommand{\natija}[1]{\cellcolor{brandOrange!15}\textbf{#1}}
% Empty cell to be filled by the student
\newcommand{\fragh}{\rule{0pt}{0.55cm}}
\newcolumntype{C}[1]{>{\centering\arraybackslash}m{#1}}
\newcolumntype{R}[1]{>{\raggedleft\arraybackslash}m{#1}}

%------------------------------------------------
% Boxes
%------------------------------------------------
\tcbset{labbox/.style={enhanced, breakable, arc=3pt, boxrule=1.2pt,
    coltitle=white, fonttitle=\khatt\bfseries,
    top=1.5mm, bottom=1.5mm, left=3mm, right=3mm,
    before skip=6pt, after skip=6pt}}
% Worked example
\newtcolorbox{mithal}[1]{labbox, colback=brandTeal!6, colframe=brandTeal,
    title={مثال محلول: #1}}
% Reminder / method
\newtcolorbox{tadhkir}[1][تذكير]{labbox, colback=brandTeal!3,
    colframe=tealDeep, title={#1}}
% Warning
\newtcolorbox{tanbih}[1][انتبه]{labbox, colback=brandOrange!8,
    colframe=brandOrange, title={#1}}
% Exercise
\newtcolorbox{tamrin}[1]{labbox, colback=white, colframe=brandDark!80,
    title={#1}}

% Section title
\newcommand{\qism}[1]{\section*{\khatt\underline{#1}}}

% Title block: \unwan{lab number}{title}{subtitle}
\newcommand{\unwan}[3]{%
    \begin{center}
        {\Large\khatt\textbf{المخبر التطبيقي رقم #1}}\\[0.15cm]
        {\Large\khatt\textbf{#2}}\\[0.1cm]
        {\khatt\textbf{#3}}
    \end{center}
    \rule{\textwidth}{1pt}
    \begin{center}
        \small المدة: ساعة ونصف \quad$\cdot$\quad الأدوات: آلة حاسبة، أداة ذكاء اصطناعي، مصدر موثوق
    \end{center}
    \rule{\textwidth}{1pt}}

% Percentage reminder used in every lab
\newcommand{\nisba}{%
    \begin{tadhkir}[تذكير: حساب النسبة المئوية]
        \begin{center}
            \textbf{النسبة المئوية = (الجزء $\div$ الكل) $\times$ 100}
        \end{center}
        مثال: 15 إجابة صحيحة من 20 سؤالًا:
        $\dfrac{15}{20}\times 100 = 75\%$.
    \end{tadhkir}}


\begin{document}

\unwan{4}{الأخلاقيات في الذكاء الاصطناعي}{المسؤولية ــ الإفصاح ــ الأمانة العلمية بالأرقام}

\qism{أهداف المخبر}
\begin{itemize}
    \item تحليل نتائج استبيان حول استعمال الطلبة للذكاء الاصطناعي.
    \item تصنيف الاستعمالات: مقبول، مقبول بشروط، غير مسؤول، مرفوض.
    \item التحقق من المراجع التي تقترحها الأداة وحساب نسبة المراجع الوهمية.
    \item حساب نسبة النقل غير الموثق في بحث، والحكم عليه.
    \item كتابة جملة إفصاح دقيقة.
\end{itemize}

\begin{tanbih}[ملاحظة حول البيانات]
    الأرقام الواردة في هذا المخبر أمثلة تعليمية، وعتبة الـ 10\% المستعملة
    في التمرين 3 افتراضية للتدريب فقط؛ كل مؤسسة تحدد قواعدها الخاصة.
\end{tanbih}

\nisba

%================================================
\qism{الجزء الأول: أمثلة محلولة}
%================================================

\begin{mithal}{استبيان: كيف يستعمل الطلبة الأداة؟}
سأل أستاذ \textbf{50 طالبًا} عن الاستعمال الرئيسي للذكاء الاصطناعي في واجباتهم
(إجابة واحدة لكل طالب):

\begin{center}
\begin{tabular}{|R{6.6cm}|C{1.8cm}|C{1.8cm}|C{3cm}|}
    \hline
    \ras{الاستعمال} & \ras{العدد} & \ras{النسبة} & \ras{الموقف} \\ \hline
    فهم قاعدة نحوية صعبة                       & 18 & 36\% & مقبول \\ \hline
    طلب أفكار، ثم الكتابة بنفسي                 & 12 & 24\% & مقبول \\ \hline
    مسودة بالأداة، ثم مراجعة وتحقق وإفصاح       & 8  & 16\% & مقبول بشروط \\ \hline
    نقل معلومات الأداة دون تحقق                & 7  & 14\% & غير مسؤول \\ \hline
    نسخ إجابة كاملة وتقديمها باسمي              & 5  & 10\% & انتحال مرفوض \\ \hline
    \textbf{المجموع}                            & \textbf{50} & 100\% & -- \\ \hline
\end{tabular}
\end{center}

\textbf{الحساب:}
\begin{itemize}
    \item الاستعمال المقبول: $18 + 12 = 30$ طالبًا، أي $\dfrac{30}{50}\times 100 = 60\%$.
    \item الاستعمال المقبول بشروط: $16\%$.
    \item الاستعمال الإشكالي (غير مسؤول + مرفوض): $7 + 5 = 12$ طالبًا، أي \natija{$24\%$}.
\end{itemize}
\textbf{الاستنتاج:} أغلب الطلبة يستعملون الأداة بطريقة مقبولة، لكن طالبًا من كل أربعة
تقريبًا يحتاج إلى توعية بالمسؤولية والأمانة العلمية.
\end{mithal}

\begin{mithal}{المراجع الوهمية}
طلب طالب من أداة ذكاء اصطناعي \textbf{10 مراجع} لبحثه عن «التشبيه في الشعر العباسي»،
ثم بحث عن كل مرجع في فهرس المكتبة:

\begin{center}
\begin{tabular}{|R{7cm}|C{2.2cm}|C{2.2cm}|}
    \hline
    \ras{نتيجة التحقق} & \ras{العدد} & \ras{النسبة} \\ \hline
    مرجع موجود ومعلوماته صحيحة                    & 6 & 60\% \\ \hline
    مرجع موجود، لكن السنة أو الناشر خاطئ           & 2 & 20\% \\ \hline
    مرجع \textbf{غير موجود أصلًا} (وهمي)           & 2 & \natija{20\%} \\ \hline
    \textbf{المجموع}                              & \textbf{10} & 100\% \\ \hline
\end{tabular}
\end{center}

\textbf{الاستنتاج:} $40\%$ من المراجع لا يمكن نقلها كما هي.
لو وضعها الطالب في بحثه دون تحقق، فهو \textbf{المسؤول} عن الخطأ، لا الأداة.
\end{mithal}

%================================================
\qism{الجزء الثاني: تمارين تطبيقية}
%================================================

\begin{tamrin}{التمرين 1: صنّف واحسب}
سأل أستاذ \textbf{40 طالبًا} عن استعمالهم الرئيسي للأداة:

\begin{center}
\begin{tabular}{|C{0.6cm}|R{6.6cm}|C{1.4cm}|C{1.6cm}|C{2.6cm}|}
    \hline
    \ras{ر} & \ras{الاستعمال} & \ras{العدد} & \ras{النسبة} & \ras{الموقف} \\ \hline
    1 & ترجمة مصطلح، ثم التحقق منه في معجم.                    & 9 & \fragh & \\ \hline
    2 & طلب شرح مبسط لبحر شعري.                                 & 8 & \fragh & \\ \hline
    3 & تدقيق إملائي لنص كتبتُه بنفسي.                          & 6 & \fragh & \\ \hline
    4 & توليد مسودة مقال، ثم إعادة كتابتها مع الإفصاح.          & 5 & \fragh & \\ \hline
    5 & نقل تواريخ وأسماء من الأداة دون تحقق.                   & 5 & \fragh & \\ \hline
    6 & حل واجب كامل بالأداة وتسليمه باسمي.                     & 4 & \fragh & \\ \hline
    7 & استعمال الأداة في امتحان يمنع ذلك.                      & 3 & \fragh & \\ \hline
\end{tabular}
\end{center}

\begin{enumerate}
    \item أكمل عمود النسبة، ثم عمود الموقف:
          (مقبول / مقبول بشروط / غير مسؤول / مرفوض).
    \item أكمل جدول التلخيص:
    \begin{center}
    \begin{tabular}{|R{3cm}|C{2.6cm}|C{2.6cm}|}
        \hline
        \ras{الموقف} & \ras{عدد الطلبة} & \ras{النسبة} \\ \hline
        مقبول         & \fragh & \\ \hline
        مقبول بشروط   & \fragh & \\ \hline
        غير مسؤول     & \fragh & \\ \hline
        مرفوض         & \fragh & \\ \hline
        \textbf{المجموع} & 40 & 100\% \\ \hline
    \end{tabular}
    \end{center}
    \item قارن نسبة الاستعمال المرفوض هنا بنسبته في المثال المحلول الأول.
\end{enumerate}
\end{tamrin}

\begin{tamrin}{التمرين 2: تحقق من المراجع}
اقترحت أداة \textbf{12 مرجعًا}. بعد التحقق في المكتبة وجد الطالب:
\textbf{5} مراجع صحيحة، و\textbf{3} مراجع موجودة بمعلومات خاطئة، و\textbf{4} مراجع غير موجودة.

\begin{enumerate}
    \item أنشئ جدولًا مثل جدول المثال المحلول الثاني، واحسب النسب.
    \item كم مرجعًا يمكن استعماله مباشرة؟ وكم مرجعًا يحتاج إلى تصحيح؟
    \item لو تحقق الطالب من 4 مراجع فقط واكتفى بها، هل يكون قد تصرف بمسؤولية؟ علّل.
    \item \textbf{تطبيق:} اطلب من أداة 5 مراجع في موضوع من اختيارك، وتحقق منها،
          ثم احسب نسبة المراجع الوهمية.
\end{enumerate}
\end{tamrin}

\begin{tamrin}{التمرين 3: احكم على البحوث}
فحص أستاذ أربعة بحوث. «النقل غير الموثق» هو كل نص منقول (من كتاب أو موقع أو أداة ذكاء
اصطناعي) دون ذكر مصدره. القاعدة الافتراضية: \textbf{النقل غير الموثق لا يتجاوز 10\%}.

\begin{center}
\begin{tabular}{|C{1.2cm}|C{1.8cm}|C{1.7cm}|C{2.1cm}|C{1.3cm}|C{2cm}|C{2cm}|}
    \hline
    \ras{البحث} & \ras{عدد الكلمات} & \ras{نقل موثق} & \ras{نقل غير موثق}
        & \ras{إفصاح} & \ras{نسبة غير الموثق} & \ras{الحكم} \\ \hline
    1 & 2000 & 300 & 0    & نعم & \fragh & \\ \hline
    2 & 2500 & 200 & 1500 & لا  & \fragh & \\ \hline
    3 & 1800 & 0   & 90   & نعم & \fragh & \\ \hline
    4 & 3000 & 450 & 600  & لا  & \fragh & \\ \hline
\end{tabular}
\end{center}

\begin{enumerate}
    \item احسب نسبة النقل غير الموثق في كل بحث
          (النقل غير الموثق $\div$ عدد الكلمات $\times$ 100).
    \item احكم على كل بحث: مقبول / مقبول مع تصحيح / انتحال.
    \item البحث 1 فيه 300 كلمة منقولة. لماذا لا يُعدّ انتحالًا؟
\end{enumerate}

\begin{tanbih}
    أدوات «كشف النص المولّد بالذكاء الاصطناعي» تعطي \textbf{تقديرًا} قد يخطئ،
    وليست دليلًا قاطعًا. الحكم النهائي يعتمد على الحوار مع الطالب ومراجعة عمله.
\end{tanbih}
\end{tamrin}

\begin{tamrin}{التمرين 4: اكتب جملة إفصاح دقيقة}
أنجزت طالبة بحثًا من خمس مراحل، واستعملت الأداة في بعضها:

\begin{center}
\begin{tabular}{|C{0.6cm}|R{4.5cm}|C{2.6cm}|R{5.2cm}|}
    \hline
    \ras{ر} & \ras{المرحلة} & \ras{استعمال الأداة} & \ras{كيف؟} \\ \hline
    1 & اختيار الموضوع       & نعم & اقتراح 5 موضوعات، اختارت واحدًا. \\ \hline
    2 & جمع المصادر          & لا  & -- \\ \hline
    3 & وضع الخطة            & نعم & خطة أولية عدّلتها. \\ \hline
    4 & الكتابة              & لا  & -- \\ \hline
    5 & التدقيق الإملائي     & نعم & تصحيح الأخطاء الإملائية فقط. \\ \hline
\end{tabular}
\end{center}

\begin{enumerate}
    \item في كم مرحلة استعملت الأداة؟ احسب النسبة من مجموع المراحل.
    \item اكتب جملة إفصاح (3 أسطر على الأكثر) تذكر فيها: اسم الأداة، والمراحل،
          وما فعلته الطالبة بنفسها.
    \item لماذا لا يكفي أن تكتب: «استعملت الذكاء الاصطناعي»؟
\end{enumerate}
\end{tamrin}

\begin{tamrin}{التمرين 5: شبكة التقييم}
يقيّم أستاذ البحوث من 20 بهذه الشبكة:

\begin{center}
\begin{tabular}{|R{4.4cm}|C{2.2cm}|C{2.4cm}|C{2.4cm}|}
    \hline
    \ras{المعيار} & \ras{العلامة القصوى} & \ras{الطالب أ} & \ras{الطالب ب} \\ \hline
    المحتوى                 & 8  & 7 & 8 \\ \hline
    المنهجية                & 4  & 3 & 4 \\ \hline
    التوثيق والمراجع         & 4  & 4 & 1 \\ \hline
    الإفصاح والأمانة العلمية & 4  & 4 & 0 \\ \hline
    \textbf{المجموع}        & \textbf{20} & \fragh & \fragh \\ \hline
\end{tabular}
\end{center}

\begin{enumerate}
    \item احسب مجموع كل طالب.
    \item الطالب ب له محتوى أفضل، فلماذا حصل على علامة أقل؟
    \item كم نقطة خسر الطالب ب بسبب التوثيق والأمانة؟ وما نسبتها من 20؟
\end{enumerate}
\end{tamrin}

%================================================
\newpage
\qism{الحلول المختصرة}
%================================================

\textbf{التمرين 1:}
النسب: $22.5\%$، $20\%$، $15\%$، $12.5\%$، $12.5\%$، $10\%$، $7.5\%$.\\
الموقف: 1، 2، 3 مقبول \quad 4 مقبول بشروط \quad 5 غير مسؤول \quad 6، 7 مرفوض.\\
التلخيص: مقبول $9+8+6 = 23$ ($57.5\%$) \quad مقبول بشروط 5 ($12.5\%$) \quad
غير مسؤول 5 ($12.5\%$) \quad مرفوض $4+3 = 7$ ($17.5\%$).\\
الاستعمال المرفوض هنا $17.5\%$ مقابل $10\%$ في المثال الأول، فهو أعلى.

\textbf{التمرين 2:}
صحيحة $\dfrac{5}{12} \approx 41.7\%$، معلومات خاطئة $\dfrac{3}{12} = 25\%$،
وهمية $\dfrac{4}{12} \approx 33.3\%$.
يستعمل 5 مباشرة، ويصحح 3، ويحذف 4.
لا: كل مرجع يُذكر في البحث يجب التحقق منه، فالطالب مسؤول عن كل مرجع يضعه.

\textbf{التمرين 3:}
النسب: البحث 1: $0\%$ \quad البحث 2: $60\%$ \quad البحث 3: $5\%$ \quad البحث 4: $20\%$.\\
الحكم: 1 مقبول \quad 2 انتحال \quad 3 مقبول مع تصحيح (توثيق الـ 90 كلمة) \quad 4 انتحال.\\
البحث 1 ليس انتحالًا لأن كل النقل \textbf{موثق} بمصادره، مع إفصاح.

\textbf{التمرين 4:}
3 مراحل من 5، أي $60\%$.
مثال: «استعملتُ أداة \textenglish{ChatGPT} لاقتراح موضوعات البحث، ولوضع خطة أولية عدّلتُها،
وللتدقيق الإملائي. جمعتُ المصادر وكتبتُ النص بنفسي، وتحققت من كل المعلومات.»
العبارة العامة لا تبيّن حجم الاستعمال ولا ما أنجزته الطالبة بنفسها، فهي غير شفافة.

\textbf{التمرين 5:}
أ: $7+3+4+4 = 18$ \quad ب: $8+4+1+0 = 13$.
خسر ب في التوثيق والأمانة: $(4-1) + (4-0) = 7$ نقاط، أي $\dfrac{7}{20}\times 100 = 35\%$.
المحتوى الجيد لا يعوّض غياب الأمانة العلمية.

\end{document}
